
\textbf{Lemme.} \uline{Soit} \(X'\xrightarrow{f}X\) \uline{un morphisme projectif et parfait}, \uline{avec} \(X\) \uline{quasi-compact}, \uline{et soit} \(\underline{L}^\bullet\) \uline{un complexe parfait sur} \(\underline{O}_{X'}\), \uline{acyclique en degrés} \(>0\). \uline{Alors il existe un entier} \(n_0\) \uline{tel que pour} \(n\ge n_0\) \uline{le complexe} \(\mathbb Rf_*\bigl(\underline{L}^\bullet\otimes_{\underline{O}_{X'}}\underline{O}_{X'}(n)\bigr)\) \uline{soit acyclique en degrés} \(>0\).

On applique ce lemme au complexe réduit à \(\underline{O}_{X'}\); il existe donc un entier \(n_0\) tel que pour \(n\ge n_0\), on ait \(R^if_*(\underline{O}_{X'}(n))=0\) si \(i\ge1\). On redescend alors par récurrence jusqu'à \(n=0\). Pour cela, on considère la suite exacte:
\[
0\longrightarrow
\underline{O}_{X'}(n+1)
\longrightarrow
\underline{O}_{X'}(n)
\longrightarrow
\underline{O}_{Y'}(n)
\longrightarrow0,
\]
et on lui applique le foncteur \(f_*\). Or on a \(\mathbb Rf_*(\underline{O}_{Y'}(n))=i_*\circ\mathbb Rg_*(\underline{O}_{Y'}(n))\). Comme \(Y'=\bfP(\underline{N})\), \(\underline{O}_{Y'}(n)\) est \(g_*\)-acyclique pour \(n\ge0\). On a donc pour tout \(n\ge0\), et tout \(i\ge1\):
\[
R^if_*(\underline{O}_{Y'}(n))=0,
\]
et par conséquent les suites exactes:
\[
R^if_*(\underline{O}_{X'}(n+1))
\longrightarrow
R^if_*(\underline{O}_{X'}(n))
\longrightarrow0,
\]
ce qui permet d'achever la récurrence.

ii) Il s'agit de montrer que le morphisme canonique:
\[
\bigoplus_{n\ge0}\underline{J}^n
\longrightarrow
\bigoplus_{n\ge0}f_*\bigl(\underline{O}_{X'}(n)\bigr)
\]
est un isomorphisme. La propriété étant locale, on peut supposer \(\underline{J}\) défini par une suite \(f_1,\ldots,f_d\) de sections de \(\underline{O}_X\), telles que l'homomorphisme \((\underline{O}_X)^d\longrightarrow\underline{O}_X\) correspondant soit régulier. Posons: \(\underline{S}=\bigoplus_{n\ge0}\underline{J}^n\); on définit un homomorphisme d'algèbres: \(\underline{O}_X[T_1,\ldots,T_d]\xrightarrow{v}\underline{S}\), surjectif, en envoyant 
